
\documentclass{article}
\usepackage{colortbl}
\usepackage{makecell}
\usepackage{multirow}
\usepackage{supertabular}

\begin{document}

\newcounter{utterance}

\centering \large Interaction Transcript for game `cladder', experiment `full\_v1.5\_default', episode 6304 with Qwen3.5{-}4B.
\vspace{24pt}

{ \footnotesize  \setcounter{utterance}{1}
\setlength{\tabcolsep}{0pt}
\begin{supertabular}{c@{$\;$}|p{.15\linewidth}@{}p{.15\linewidth}p{.15\linewidth}p{.15\linewidth}p{.15\linewidth}p{.15\linewidth}}
    \# & \multicolumn{2}{c}{Player} && \multicolumn{2}{c}{Game Master} \\
    \hline

    \theutterance \stepcounter{utterance}  
    & & & \multicolumn{4}{p{0.6\linewidth}}{
        \cellcolor[rgb]{0.9,0.9,0.9}{
            \makecell[{{p{\linewidth}}}]{
                \texttt{\tiny{[P1$\langle$GM]}}
                \texttt{You are an expert in causal inference. The following question is not a typical commonsense query, but rather a meticulously designed question created by a professor specializing in causal inference, intended to assess the students' mastery of the course content.} \\
\texttt{Imagine a self{-}contained, hypothetical world with only the following conditions, and without any unmentioned factors or causal relationships: Zuph has a direct effect on rixq and xevu. Rixq has a direct effect on xevu. The overall probability of rixq is 32\%. The probability of not rixq and xevu is 35\%. The probability of rixq and xevu is 21\%. Is the chance of xevu smaller when observing rixq?} \\
\texttt{Start your answer with 'Yes' or 'No', followed by additional reasoning or evidence to support your explanation.} \\
            }
        }
    }
     \\ \\

    \theutterance \stepcounter{utterance}  
    & \multicolumn{4}{p{0.6\linewidth}}{
        \cellcolor[rgb]{0.561,0.737,0.561}{
            \makecell[{{p{\linewidth}}}]{
                \texttt{\tiny{[P1$\rangle$GM]}}
                \texttt{Yes} \\
\\ 
\texttt{**Reasoning:**} \\
\\ 
\texttt{To determine whether the chance of \$xevu\$ is smaller when observing \$rixq\$, we must compare the marginal probability of \$xevu\$ (\$P(xevu)\$) with the conditional probability of \$xevu\$ given \$rixq\$ (\$P(xevu | rixq)\$).} \\
\\ 
\texttt{1.  **Calculate the marginal probability of \$xevu\$ (\$P(xevu)\$):**} \\
\texttt{    We are given the joint probabilities for the combinations of \$rixq\$ and \$xevu\$. Let \$R\$ represent \$rixq\$ and \$X\$ represent \$xevu\$.} \\
\texttt{    *   \$P(R) = 0.32\$} \\
\texttt{    *   \$P(\textbackslash{}neg R \textbackslash{}cap X) = 0.35\$} \\
\texttt{    *   \$P(R \textbackslash{}cap X) = 0.21\$} \\
\\ 
\texttt{    The total probability of \$X\$ is the sum of the probabilities of \$X\$ occurring with \$R\$ and \$X\$ occurring without \$R\$:} \\
\texttt{    \$\$P(X) = P(R \textbackslash{}cap X) + P(\textbackslash{}neg R \textbackslash{}cap X)\$\$} \\
\texttt{    \$\$P(X) = 0.21 + 0.35 = 0.56\$\$} \\
\\ 
\texttt{2.  **Calculate the conditional probability of \$xevu\$ given \$rixq\$ (\$P(X | R)\$):**} \\
\texttt{    Using the definition of conditional probability:} \\
\texttt{    \$\$P(X | R) = \textbackslash{}frac\{P(R \textbackslash{}cap X)\}\{P(R)\}\$\$} \\
\texttt{    \$\$P(X | R) = \textbackslash{}frac\{0.21\}\{0.32\} \textbackslash{}approx 0.65625\$\$} \\
\\ 
\texttt{3.  **Compare the probabilities:**} \\
\texttt{    *   Marginal probability \$P(xevu) = 0.56\$} \\
\texttt{    *   Conditional probability \$P(xevu | rixq) \textbackslash{}approx 0.656\$} \\
\\ 
\texttt{    Since \$0.656 > 0.56\$, the chance of \$xevu\$ is actually **larger** (not smaller) when observing \$rixq\$.} \\
\\ 
\texttt{**Conclusion:**} \\
\texttt{The observation of \$rixq\$ increases the probability of \$xevu\$ from 56\% to approximately 65.6\%. Therefore, the answer to the question "Is the chance of xevu smaller when observing rixq?" is **No**.} \\
            }
        }
    }
    & & \\ \\

    \theutterance \stepcounter{utterance}  
    & & & \multicolumn{2}{p{0.3\linewidth}}{
        \cellcolor[rgb]{0.9,0.9,0.9}{
            \makecell[{{p{\linewidth}}}]{
                \texttt{\tiny{[GM$|$GM]}}
                \texttt{Yes} \\
            }
        }
    }
    & & \\ \\

    \theutterance \stepcounter{utterance}  
    & & & \multicolumn{2}{p{0.3\linewidth}}{
        \cellcolor[rgb]{0.9,0.9,0.9}{
            \makecell[{{p{\linewidth}}}]{
                \texttt{\tiny{[GM$|$GM]}}
                \texttt{no} \\
            }
        }
    }
    & & \\ \\

    \theutterance \stepcounter{utterance}  
    & & & \multicolumn{2}{p{0.3\linewidth}}{
        \cellcolor[rgb]{0.9,0.9,0.9}{
            \makecell[{{p{\linewidth}}}]{
                \texttt{\tiny{[GM$|$GM]}}
                \texttt{game\_result = LOSE} \\
            }
        }
    }
    & & \\ \\

\end{supertabular}
}

\end{document}
